\documentclass{article}
\usepackage{amsmath,amssymb}
\usepackage{xurl}
\usepackage[hidelinks]{hyperref}
% Shared commands for notes in docs/paper-gaps/.

\DeclareUnicodeCharacter{2080}{\ifmmode{}_{0}\else\textsubscript{0}\fi}
\DeclareUnicodeCharacter{2081}{\ifmmode{}_{1}\else\textsubscript{1}\fi}
\DeclareUnicodeCharacter{2082}{\ifmmode{}_{2}\else\textsubscript{2}\fi}
\DeclareUnicodeCharacter{2099}{\ifmmode{}_{n}\else\textsubscript{n}\fi}

\newcommand{\Lean}{\textsc{Lean}}
\ExplSyntaxOn
\NewDocumentCommand{\leanid}{m}{
  \ifmmode
    \textsf{\detokenize{#1}}
  \else
    \group_begin:
    \urlstyle{sf}
    \tl_set:Nn \l_tmpa_tl {#1}
    \tl_replace_all:Nnn \l_tmpa_tl {\_} {_}
    \exp_args:NV \path \l_tmpa_tl
    \group_end:
  \fi
}
\ExplSyntaxOff
\providecommand{\tr}{\operatorname{tr}}
\newcommand{\tnleanIssueBase}{https://github.com/LionSR/TNLean/issues/}
\newcommand{\tnleanPrBase}{https://github.com/LionSR/TNLean/pull/}
\newcommand{\tnleanDocBase}{https://sirui-lu.com/TNLean/docs/}
\newcommand{\ghissue}[1]{\href{\tnleanIssueBase#1}{GitHub issue \##1}}
\newcommand{\ghpr}[1]{\href{\tnleanPrBase#1}{PR \##1}}
\newcommand{\doclink}[1]{\href{\tnleanDocBase#1.html}{\path{#1}}}

% Dirac notation and the identity operator, as in blueprint/src/macros/common.tex.
% \providecommand keeps note-local definitions authoritative.
\usepackage{dsfont}
\providecommand{\ket}[1]{|#1\rangle}
\providecommand{\bra}[1]{\langle #1|}
\providecommand{\braket}[2]{\langle #1 | #2 \rangle}
\providecommand{\ketbra}[2]{|#1\rangle\!\langle #2|}
\providecommand{\Id}{\mathds{1}}
\providecommand{\one}{\mathds{1}}
\providecommand{\C}{\mathbb{C}}
\providecommand{\MN}[1]{M_{#1}(\C)}
\providecommand{\MD}{M_D(\C)}

% Verdict marker: \gapnote{<kind>}{<status>}. Kind is the policy
% classification (clarification, local-correction, scope-restriction,
% unfaithful, false-source, open-gap); status is open, wip (elimination
% underway), resolved, or historical. Machine-read by the site index;
% prints a verdict line.
\newcommand{\gapnote}[2]{\par\noindent\textbf{Verdict:} #1 (#2).\par}

\title{Inhomogeneous Matrix Product States: Common Bond Dimension and Injective Blocks}
\author{}
\date{2026-10-01}
\begin{document}
\maketitle
\gapnote{scope-restriction}{open}

\section{What the source states}
The paragraph ``Inhomogeneous short-range correlated MPS''
of~\cite{Malz2024Preparation} considers a sequence of matrix product states
$\ket{\phi_N}$ ``with bond dimension at most $D$'' that are not translation
invariant. The sequence has finite correlation length if, after blocking
$q=O(\log N)$ sites, the state $\ket{\phi_{\mathrm{pos}}}$ of the positive parts
$P_i$ of the polar decompositions of the blocked tensors (eq.~(18)) is
approximated by a product of nearest-neighbour pairs
$\ket\Omega=\bigotimes_{i=1}^{N/q}\ket{\omega^i}_{R_iL_{i+1}}$ with
$\varepsilon(\Omega,\phi_{\mathrm{pos}})\to0$. The paragraph then asserts that
preparing $\ket\Omega$ and applying the isometries of the polar decompositions
gives depth $O(\log(N/\varepsilon))$ with error
$\varepsilon(\Omega,\phi_{\mathrm{pos}})$.

\section{The restriction adopted}
The formal statements\footnote{Formal declarations:
    \leanid{MPSPreparation.exists_isPreparedInDepth_inhomogeneous} and
    \leanid{MPSPreparation.exists_isPreparedInDepth_of_isPairApproximable}
    in \doclink{TNLean/MPS/Preparation/InhomogeneousPreparation}.}
work on a ring of $N$ sites cut into blocks of lengths $\ell_k$ and assume three
things.
\begin{enumerate}
    \item Every bond of the ring has the same dimension $D$. The source allows
          bond dimensions at most $D$, which may vary along the ring.
    \item Every blocked tensor is injective, so that each isometric factor
          $V_k$ satisfies $V_k^\dagger V_k=\Id_{D^2}$. The source makes this
          assumption for its polar decompositions in the footnote to the
          paragraph ``Approximation through the fixed-point state'' and uses it
          silently here: the identity between the error of
          $(\bigotimes_kV_k)\ket\Omega$ against $\ket{\phi_N}$ and
          $\varepsilon(\Omega,\phi_{\mathrm{pos}})$ is proved when the $V_k$
          are isometries; for partial isometries the projections $\Pi_k$
          enter.
    \item Every block has at least $3D$ sites, which the circuit needs to encode
          the two bond indices of a block and run the staircase between them.
          For $q=O(\log N)$ growing with $N$ this holds for large $N$.
\end{enumerate}
The finite-correlation assumption is stated for one ring: some unit pairs
$\omega^k$ have $\varepsilon(\Omega,\phi_{\mathrm{pos}})\le\delta$. The formal
statements prove depth at most $CL$ for blocks of lengths at most $L$, with $C$
depending only on $d$ and $D$, and error exactly
$\varepsilon(\Omega,\phi_{\mathrm{pos}})$, hence at most $\delta$. The source's
total depth $O(\log(N/\varepsilon))$ is the choice $L=O(\log(N/\varepsilon))$,
which is not made here; the source's definition fixes blocks of length
$q=O(\log N)$.

\section{Nonzero state of the positive parts}
The error $\varepsilon(\Omega,\phi_{\mathrm{pos}})$ compares $\ket\Omega$ with
the normalized state $\ket{\phi_{\mathrm{pos}}}/\lVert\phi_{\mathrm{pos}}\rVert$,
which presupposes $\ket{\phi_{\mathrm{pos}}}\neq0$. The formal approximation
condition\footnote{Formal declaration:
    \leanid{MPSPreparation.IsPairApproximable}.}
states $\ket{\phi_{\mathrm{pos}}}\neq0$ explicitly; without it the normalized
vector would be read as $0$ and the condition would hold for every
$\delta\ge1$. This is a local fix of the intended reading, not a restriction.

\section{Elimination}
\begin{enumerate}
    \item \emph{Varying bond dimensions.} Replace the chain of square tensors by
          one with bond dimensions $D_i\le D$. The blocked tensor of block $k$
          is then a map $\mathbb C^{D_{i}}\otimes\mathbb C^{D_{j}}\to
              (\mathbb C^d)^{\otimes\ell_k}$, and the pair $\omega^k$ lives on
          $\mathbb C^{D_{j}}\otimes\mathbb C^{D_{j}}$; the encoding of the bond
          indices in $D$ sites and the sequential factorization go through
          unchanged with rectangular matrices.
    \item \emph{Injectivity.} Without it $V_k$ is a partial isometry with
          $V_k^\dagger V_k=\Pi_k$, and the error becomes that of
          $(\bigotimes_k\Pi_k)\ket\Omega$ against $\ket{\phi_{\mathrm{pos}}}$,
          as in the non-normal case of the source.
    \item \emph{Block length.} The bound $3D\le\ell_k$ is needed only by
          the circuit, not by the state. The upper bound $q=O(\log N)$ alone
          allows bounded block lengths, so the sequence-level statement reads
          the source's blocking as $q_N\ge c\log N$ for some $c>0$, a block
          length growing with $N$ as in the translation-invariant case. Then
          $q_N\ge3D$ once $N\ge N_0$, with $N_0$ depending only on $D$ and
          $c$; each of the finitely many rings with $N<N_0$ is prepared exactly
          by a circuit of depth bounded in terms of $N_0$ and $d$, and the
          largest of these depths is absorbed into the constant of the depth
          bound.
    \item \emph{Sequence-level statement.} State the source's definition of
          finite correlation length for the sequence, with a rate: unit pairs
          with $\varepsilon(\Omega,\phi_{\mathrm{pos}})\le\delta(N,q)$ after
          blocking $q$ sites. Choosing the block length $L$ with
          $\delta(N,L)\le\varepsilon$, which is $L=O(\log(N/\varepsilon))$
          when $\delta(N,q)$ decays exponentially in $q$ at fixed $N$ up to a
          factor polynomial in $N$, as in the translation-invariant case, turns
          the depth bound $CL$ and the error $\delta$ proved for one ring into
          the source's depth $O(\log(N/\varepsilon))$ with error at most
          $\varepsilon$.
\end{enumerate}

\bibliographystyle{alpha}
\bibliography{references}
\end{document}
